%
\begin{isabellebody}%
\setisabellecontext{Week{\isacharunderscore}{\kern0pt}{\isadigit{1}}{\isacharunderscore}{\kern0pt}sol}%
%
\isadelimtheory
\isanewline
%
\endisadelimtheory
%
\isatagtheory
\isacommand{theory}\isamarkupfalse%
\ Week{\isacharunderscore}{\kern0pt}{\isadigit{1}}{\isacharunderscore}{\kern0pt}sol\isanewline
\isakeyword{imports}\ Main\isanewline
\isakeyword{begin}%
\endisatagtheory
{\isafoldtheory}%
%
\isadelimtheory
%
\endisadelimtheory
%
\begin{isamarkuptext}%
\vspace{15ex}\ExerciseSheet{6}{}%
\end{isamarkuptext}\isamarkuptrue%
%
\begin{isamarkuptext}%
Before you start, please open a new Isabelle/HOL theory file called Week\_1.thy and add the
       following to its beginning.

\noindent theory Week\_1\\
imports Main\\
begin%
\end{isamarkuptext}\isamarkuptrue%
%
\begin{isamarkuptext}%
\Exercise{Addition Commutes}%
\end{isamarkuptext}\isamarkuptrue%
%
\begin{isamarkuptext}%
Consider the function \isa{{\isacharparenleft}{\kern0pt}{\isacharplus}{\kern0pt}{\isacharparenright}{\kern0pt}} defined in Isabelle/HOL for natural numbers.

Prove the following theorem, first pen-and-paper, and then formally in Isabelle.%
\end{isamarkuptext}\isamarkuptrue%
\isacommand{lemma}\isamarkupfalse%
\ add{\isacharunderscore}{\kern0pt}commutes{\isacharcolon}{\kern0pt}\ {\isachardoublequoteopen}{\isacharparenleft}{\kern0pt}n{\isacharcolon}{\kern0pt}{\isacharcolon}{\kern0pt}nat{\isacharparenright}{\kern0pt}\ {\isacharplus}{\kern0pt}\ m\ {\isacharequal}{\kern0pt}\ m\ {\isacharplus}{\kern0pt}\ n{\isachardoublequoteclose}%
\begin{isamarkuptext}%
\paragraph{Hints:} 

 1. Perform the proof by inuction on \isa{n}, using substitutions \isa{subst}, and forward reasoning using
   the method \isa{rule}.

 2. You will need to identify multiple lemmas about how \isa{add} works. Do not prove those lemmas. Use 
    \isa{sorry} as the proof. This is the top-down approach. 

 3. There is a proof which uses the following lemmas:

   \isa{add{\isacharunderscore}{\kern0pt}{\isadigit{0}}{\isacharunderscore}{\kern0pt}right{\isacharcomma}{\kern0pt}\ add{\isacharunderscore}{\kern0pt}{\isadigit{0}}{\isacharunderscore}{\kern0pt}left{\isacharcomma}{\kern0pt}\ add{\isacharunderscore}{\kern0pt}Suc{\isacharunderscore}{\kern0pt}right{\isacharcomma}{\kern0pt}\ add{\isacharunderscore}{\kern0pt}Suc}%
\end{isamarkuptext}\isamarkuptrue%
%
\isadelimproof
%
\endisadelimproof
%
\isatagproof
\isacommand{proof}\isamarkupfalse%
{\isacharparenleft}{\kern0pt}induction\ m{\isacharparenright}{\kern0pt}\isanewline
\ \ \isacommand{case}\isamarkupfalse%
\ {\isadigit{0}}\isanewline
\ \ \isacommand{have}\isamarkupfalse%
\ {\isadigit{2}}{\isacharcolon}{\kern0pt}\ {\isachardoublequoteopen}n\ {\isacharplus}{\kern0pt}\ {\isadigit{0}}\ {\isacharequal}{\kern0pt}\ n{\isachardoublequoteclose}\isanewline
\ \ \ \ \isacommand{apply}\isamarkupfalse%
{\isacharparenleft}{\kern0pt}subst\ {\isacharparenleft}{\kern0pt}{\isadigit{1}}{\isacharparenright}{\kern0pt}\ add{\isacharunderscore}{\kern0pt}{\isadigit{0}}{\isacharunderscore}{\kern0pt}right{\isacharparenright}{\kern0pt}\isanewline
\ \ \ \ \isacommand{{\isachardot}{\kern0pt}{\isachardot}{\kern0pt}}\isamarkupfalse%
\isanewline
\ \ \isacommand{have}\isamarkupfalse%
\ {\isadigit{4}}{\isacharcolon}{\kern0pt}\ {\isachardoublequoteopen}{\isadigit{0}}\ {\isacharplus}{\kern0pt}\ n\ {\isacharequal}{\kern0pt}\ n{\isachardoublequoteclose}\isanewline
\ \ \ \ \isacommand{apply}\isamarkupfalse%
{\isacharparenleft}{\kern0pt}subst\ {\isacharparenleft}{\kern0pt}{\isadigit{1}}{\isacharparenright}{\kern0pt}\ add{\isacharunderscore}{\kern0pt}{\isadigit{0}}{\isacharunderscore}{\kern0pt}left{\isacharparenright}{\kern0pt}\isanewline
\ \ \ \ \isacommand{{\isachardot}{\kern0pt}{\isachardot}{\kern0pt}}\isamarkupfalse%
\isanewline
\isanewline
\ \ \isacommand{have}\isamarkupfalse%
\ {\isadigit{5}}{\isacharcolon}{\kern0pt}\ {\isachardoublequoteopen}n\ {\isacharplus}{\kern0pt}\ {\isadigit{0}}\ {\isacharequal}{\kern0pt}\ n{\isachardoublequoteclose}\isanewline
\ \ \ \ \isacommand{using}\isamarkupfalse%
\ {\isadigit{2}}\isanewline
\ \ \ \ \isacommand{{\isachardot}{\kern0pt}}\isamarkupfalse%
\isanewline
\ \ \isacommand{from}\isamarkupfalse%
\ {\isadigit{5}}\ \isacommand{show}\isamarkupfalse%
\ {\isacharquery}{\kern0pt}case\isanewline
\ \ \ \ \isacommand{by}\isamarkupfalse%
{\isacharparenleft}{\kern0pt}subst\ {\isadigit{4}}{\isacharparenright}{\kern0pt}\isanewline
\isacommand{next}\isamarkupfalse%
\isanewline
\ \ \isacommand{case}\isamarkupfalse%
\ {\isacharparenleft}{\kern0pt}Suc\ m{\isacharparenright}{\kern0pt}\isanewline
\isanewline
\ \ \isacommand{hence}\isamarkupfalse%
\ {\isadigit{2}}{\isacharcolon}{\kern0pt}\ {\isachardoublequoteopen}n\ {\isacharplus}{\kern0pt}\ {\isacharparenleft}{\kern0pt}Suc\ m{\isacharparenright}{\kern0pt}\ {\isacharequal}{\kern0pt}\ Suc\ {\isacharparenleft}{\kern0pt}n\ {\isacharplus}{\kern0pt}\ m{\isacharparenright}{\kern0pt}{\isachardoublequoteclose}\isanewline
\ \ \ \ \isacommand{apply}\isamarkupfalse%
{\isacharparenleft}{\kern0pt}subst\ add{\isacharunderscore}{\kern0pt}Suc{\isacharunderscore}{\kern0pt}right{\isacharparenright}{\kern0pt}\isanewline
\ \ \ \ \isacommand{{\isachardot}{\kern0pt}{\isachardot}{\kern0pt}}\isamarkupfalse%
\isanewline
\isanewline
\ \ \isacommand{hence}\isamarkupfalse%
\ {\isadigit{4}}{\isacharcolon}{\kern0pt}\ {\isachardoublequoteopen}{\isacharparenleft}{\kern0pt}Suc\ m{\isacharparenright}{\kern0pt}\ {\isacharplus}{\kern0pt}\ n\ {\isacharequal}{\kern0pt}\ Suc\ {\isacharparenleft}{\kern0pt}m\ {\isacharplus}{\kern0pt}\ n{\isacharparenright}{\kern0pt}{\isachardoublequoteclose}\isanewline
\ \ \ \ \isacommand{apply}\isamarkupfalse%
{\isacharparenleft}{\kern0pt}subst\ add{\isacharunderscore}{\kern0pt}Suc{\isacharparenright}{\kern0pt}\isanewline
\ \ \ \ \isacommand{{\isachardot}{\kern0pt}{\isachardot}{\kern0pt}}\isamarkupfalse%
\isanewline
\isanewline
\ \ \isacommand{hence}\isamarkupfalse%
\ {\isadigit{6}}{\isacharcolon}{\kern0pt}\ {\isachardoublequoteopen}Suc\ {\isacharparenleft}{\kern0pt}n\ {\isacharplus}{\kern0pt}\ m{\isacharparenright}{\kern0pt}\ {\isacharequal}{\kern0pt}\ Suc\ {\isacharparenleft}{\kern0pt}m\ {\isacharplus}{\kern0pt}\ n{\isacharparenright}{\kern0pt}{\isachardoublequoteclose}\isanewline
\ \ \ \ \isanewline
\ \ \ \ \isacommand{thm}\isamarkupfalse%
\ Suc{\isachardot}{\kern0pt}IH\isanewline
\ \ \ \ \isanewline
\ \ \ \ \isacommand{apply}\isamarkupfalse%
{\isacharparenleft}{\kern0pt}subst\ Suc{\isachardot}{\kern0pt}IH{\isacharparenright}{\kern0pt}\isanewline
\ \ \ \ \isacommand{{\isachardot}{\kern0pt}{\isachardot}{\kern0pt}}\isamarkupfalse%
\isanewline
\isanewline
\ \ \isacommand{hence}\isamarkupfalse%
\ {\isadigit{7}}{\isacharcolon}{\kern0pt}\ {\isachardoublequoteopen}n\ {\isacharplus}{\kern0pt}\ {\isacharparenleft}{\kern0pt}Suc\ m{\isacharparenright}{\kern0pt}\ {\isacharequal}{\kern0pt}\ Suc\ {\isacharparenleft}{\kern0pt}m\ {\isacharplus}{\kern0pt}\ n{\isacharparenright}{\kern0pt}{\isachardoublequoteclose}\isanewline
\ \ \ \ \isacommand{by}\isamarkupfalse%
{\isacharparenleft}{\kern0pt}subst\ {\isacharparenleft}{\kern0pt}asm{\isacharparenright}{\kern0pt}\ {\isadigit{2}}{\isacharbrackleft}{\kern0pt}symmetric{\isacharbrackright}{\kern0pt}{\isacharparenright}{\kern0pt}\isanewline
\isanewline
\ \ \isacommand{thus}\isamarkupfalse%
\ {\isacharquery}{\kern0pt}case\isanewline
\ \ \ \ \isacommand{by}\isamarkupfalse%
{\isacharparenleft}{\kern0pt}subst\ {\isacharparenleft}{\kern0pt}asm{\isacharparenright}{\kern0pt}\ {\isadigit{4}}{\isacharbrackleft}{\kern0pt}symmetric{\isacharbrackright}{\kern0pt}{\isacharparenright}{\kern0pt}\isanewline
\isanewline
\isacommand{qed}\isamarkupfalse%
%
\endisatagproof
{\isafoldproof}%
%
\isadelimproof
%
\endisadelimproof
%
\begin{isamarkuptext}%
\Exercise{Multiplication is Monotone}%
\end{isamarkuptext}\isamarkuptrue%
%
\begin{isamarkuptext}%
Consider the multiplication function \isa{{\isacharparenleft}{\kern0pt}{\isacharasterisk}{\kern0pt}{\isacharparenright}{\kern0pt}} defined in Isabelle. Prove the following
      property of it, i.e.\ that it is monotonically increasing in both arguments. You should prove
      it both, pen-and-paper and in Isabelle.%
\end{isamarkuptext}\isamarkuptrue%
\isacommand{lemma}\isamarkupfalse%
\ mult{\isacharunderscore}{\kern0pt}mono{\isacharcolon}{\kern0pt}\isanewline
\ \ \isakeyword{fixes}\ a\ b\ c\ d{\isacharcolon}{\kern0pt}{\isacharcolon}{\kern0pt}\ nat\ %
\isamarkupcmt{Note the alternative way of fixing the types of the constants%
}\isanewline
\ \ \isakeyword{assumes}\ {\isachardoublequoteopen}a\ {\isasymle}\ b{\isachardoublequoteclose}\ {\isachardoublequoteopen}c\ {\isasymle}\ d{\isachardoublequoteclose}\isanewline
\ \ \isakeyword{shows}\ {\isachardoublequoteopen}a\ {\isacharasterisk}{\kern0pt}\ c\ {\isasymle}\ b\ {\isacharasterisk}{\kern0pt}\ d{\isachardoublequoteclose}%
\begin{isamarkuptext}%
Your pen-and-paper proof should indicate

 1. what lemmas are you assuming,

 2. on what variable are you preforming induction,

 3. what are assumptions and the proof goals in the base case and the step case, and what is the
    induction hypothesis, and

 4. how are you instantiating the induction hypothesis, i.e. how do you show that its assumptions
    hold and which quantified variables are instantiated with which constants.

\paragraph{Hints:}
 1. The proof should be by induction

  2. If the induction hypothesis is too week, you might need to strengthen by generalising/
     universally quantifying over one of the variables. E.g. if you want to universally quantify on
     a variable \isa{x}, then you can use \isa{induction\ {\isachardot}{\kern0pt}{\isachardot}{\kern0pt}\ arbitrary{\isacharcolon}{\kern0pt}\ x}.

  3. The following lemmas suffice to prove the goal with only substitution, forward reasoning, and
     induction:

     \isa{refl{\isacharcomma}{\kern0pt}\ diff{\isacharunderscore}{\kern0pt}le{\isacharunderscore}{\kern0pt}self{\isacharcomma}{\kern0pt}\ mult{\isacharunderscore}{\kern0pt}le{\isacharunderscore}{\kern0pt}mono{\isadigit{2}}{\isacharcomma}{\kern0pt}\ mult{\isacharunderscore}{\kern0pt}zero{\isacharunderscore}{\kern0pt}right{\isacharcomma}{\kern0pt}\ le{\isadigit{0}}{\isacharcomma}{\kern0pt}\ mult{\isacharunderscore}{\kern0pt}Suc{\isacharunderscore}{\kern0pt}right{\isacharcomma}{\kern0pt}\ add{\isacharunderscore}{\kern0pt}mono{\isacharcomma}{\kern0pt}\ Suc{\isacharunderscore}{\kern0pt}eq{\isacharunderscore}{\kern0pt}plus{\isadigit{1}}{\isacharcomma}{\kern0pt}\ le{\isacharunderscore}{\kern0pt}diff{\isacharunderscore}{\kern0pt}conv}

  4. When you have a term like \isa{x\ {\isacharminus}{\kern0pt}\ y}, where \isa{x} is a \isa{nat}, in the goal or the assumptions, you 
     should perform proof by case analysis on whether \isa{y\ {\isasymle}\ x}. This is because of the way \isa{{\isacharminus}{\kern0pt}} is
     defined for natural numbers, where \isa{x\ {\isacharminus}{\kern0pt}\ y\ {\isacharequal}{\kern0pt}\ {\isadigit{0}}}, for any \isa{x\ {\isasymle}\ y}.%
\end{isamarkuptext}\isamarkuptrue%
%
\isadelimproof
\ \ %
\endisadelimproof
%
\isatagproof
\isacommand{using}\isamarkupfalse%
\ assms\isanewline
\isacommand{proof}\isamarkupfalse%
{\isacharparenleft}{\kern0pt}induction\ d\ arbitrary{\isacharcolon}{\kern0pt}\ c{\isacharparenright}{\kern0pt}\isanewline
\ \ \isacommand{case}\isamarkupfalse%
\ {\isadigit{0}}\isanewline
\ \ \isacommand{then}\isamarkupfalse%
\ \isacommand{show}\isamarkupfalse%
\ {\isacharquery}{\kern0pt}case\isanewline
\ \ \ \ \isacommand{by}\isamarkupfalse%
\ auto\isanewline
\isacommand{next}\isamarkupfalse%
\isanewline
\ \ \isacommand{case}\isamarkupfalse%
\ {\isacharparenleft}{\kern0pt}Suc\ d{\isacharparenright}{\kern0pt}\isanewline
\ \ \isacommand{from}\isamarkupfalse%
\ {\isacartoucheopen}c\ {\isasymle}\ Suc\ d{\isacartoucheclose}\ \isacommand{have}\isamarkupfalse%
\ {\isadigit{1}}{\isacharcolon}{\kern0pt}\ {\isachardoublequoteopen}c\ {\isasymle}\ d\ {\isacharplus}{\kern0pt}\ {\isadigit{1}}{\isachardoublequoteclose}\isanewline
\ \ \ \ \isacommand{apply}\isamarkupfalse%
\ {\isacharparenleft}{\kern0pt}subst\ Suc{\isacharunderscore}{\kern0pt}eq{\isacharunderscore}{\kern0pt}plus{\isadigit{1}}{\isacharbrackleft}{\kern0pt}symmetric{\isacharbrackright}{\kern0pt}{\isacharparenright}{\kern0pt}\isanewline
\ \ \ \ \isacommand{{\isachardot}{\kern0pt}}\isamarkupfalse%
\isanewline
\ \ \isacommand{hence}\isamarkupfalse%
\ {\isachardoublequoteopen}c\ {\isacharminus}{\kern0pt}\ {\isadigit{1}}\ {\isasymle}\ d{\isachardoublequoteclose}\isanewline
\ \ \ \ \isacommand{apply}\isamarkupfalse%
{\isacharparenleft}{\kern0pt}subst\ le{\isacharunderscore}{\kern0pt}diff{\isacharunderscore}{\kern0pt}conv{\isacharparenright}{\kern0pt}\isanewline
\ \ \ \ \isacommand{{\isachardot}{\kern0pt}}\isamarkupfalse%
\isanewline
\isanewline
\ \ \isanewline
\isanewline
\ \ \isacommand{hence}\isamarkupfalse%
\ {\isadigit{2}}{\isacharcolon}{\kern0pt}\ {\isachardoublequoteopen}a\ {\isacharasterisk}{\kern0pt}\ {\isacharparenleft}{\kern0pt}c\ {\isacharminus}{\kern0pt}\ {\isadigit{1}}{\isacharparenright}{\kern0pt}\ {\isasymle}\ b\ {\isacharasterisk}{\kern0pt}\ d{\isachardoublequoteclose}\isanewline
\ \ \ \ \isanewline
\ \ \ \ \isacommand{apply}\isamarkupfalse%
{\isacharparenleft}{\kern0pt}rule\ Suc{\isachardot}{\kern0pt}IH{\isacharbrackleft}{\kern0pt}OF\ {\isacartoucheopen}a\ {\isasymle}\ b{\isacartoucheclose}{\isacharbrackright}{\kern0pt}{\isacharparenright}{\kern0pt}\isanewline
\ \ \ \ \isacommand{{\isachardot}{\kern0pt}}\isamarkupfalse%
\isanewline
\isanewline
\ \ \ \isanewline
\ \ \isanewline
\ \ \isacommand{show}\isamarkupfalse%
\ {\isacharquery}{\kern0pt}case\isanewline
\ \ \isacommand{proof}\isamarkupfalse%
{\isacharparenleft}{\kern0pt}cases\ c{\isacharparenright}{\kern0pt}\isanewline
\ \ \ \ \isacommand{case}\isamarkupfalse%
\ {\isadigit{0}}\isanewline
\ \ \ \ \isacommand{have}\isamarkupfalse%
\ {\isachardoublequoteopen}a\ {\isacharasterisk}{\kern0pt}\ c\ {\isacharequal}{\kern0pt}\ a\ {\isacharasterisk}{\kern0pt}\ {\isadigit{0}}{\isachardoublequoteclose}\isanewline
\ \ \ \ \ \ \isacommand{apply}\isamarkupfalse%
{\isacharparenleft}{\kern0pt}subst\ {\isadigit{0}}{\isacharparenright}{\kern0pt}\isanewline
\ \ \ \ \ \ \isacommand{{\isachardot}{\kern0pt}{\isachardot}{\kern0pt}}\isamarkupfalse%
\isanewline
\ \ \ \ \isacommand{hence}\isamarkupfalse%
\ {\isadigit{3}}{\isacharcolon}{\kern0pt}\ {\isachardoublequoteopen}a\ {\isacharasterisk}{\kern0pt}\ c\ {\isacharequal}{\kern0pt}\ {\isadigit{0}}{\isachardoublequoteclose}\isanewline
\ \ \ \ \ \ \isacommand{apply}\isamarkupfalse%
{\isacharparenleft}{\kern0pt}subst\ mult{\isacharunderscore}{\kern0pt}zero{\isacharunderscore}{\kern0pt}right{\isacharbrackleft}{\kern0pt}symmetric{\isacharbrackright}{\kern0pt}{\isacharparenright}{\kern0pt}\isanewline
\ \ \ \ \ \ \isacommand{{\isachardot}{\kern0pt}}\isamarkupfalse%
\isanewline
\ \ \ \ \isacommand{have}\isamarkupfalse%
\ {\isadigit{4}}{\isacharcolon}{\kern0pt}\ {\isachardoublequoteopen}{\isadigit{0}}\ {\isasymle}\ b\ {\isacharasterisk}{\kern0pt}\ {\isacharparenleft}{\kern0pt}Suc\ d{\isacharparenright}{\kern0pt}{\isachardoublequoteclose}\isanewline
\ \ \ \ \ \ \isacommand{using}\isamarkupfalse%
\ le{\isadigit{0}}\isanewline
\ \ \ \ \ \ \isacommand{{\isachardot}{\kern0pt}}\isamarkupfalse%
\isanewline
\ \ \ \ \isacommand{thus}\isamarkupfalse%
\ {\isacharquery}{\kern0pt}thesis\isanewline
\ \ \ \ \ \ \isacommand{apply}\isamarkupfalse%
\ {\isacharparenleft}{\kern0pt}subst\ {\isadigit{3}}{\isacharparenright}{\kern0pt}\isanewline
\ \ \ \ \ \ \isacommand{{\isachardot}{\kern0pt}}\isamarkupfalse%
\isanewline
\ \ \isacommand{next}\isamarkupfalse%
\isanewline
\ \ \ \ \isacommand{case}\isamarkupfalse%
\ {\isacharparenleft}{\kern0pt}Suc\ nat{\isacharparenright}{\kern0pt}\isanewline
\ \ \ \ \isacommand{hence}\isamarkupfalse%
\ {\isadigit{3}}{\isacharcolon}{\kern0pt}\ {\isachardoublequoteopen}a\ {\isacharasterisk}{\kern0pt}\ nat\ {\isasymle}\ b\ {\isacharasterisk}{\kern0pt}\ d{\isachardoublequoteclose}\isanewline
\ \ \ \ \ \ \isacommand{using}\isamarkupfalse%
\ {\isadigit{2}}\isanewline
\ \ \ \ \ \ \isacommand{by}\isamarkupfalse%
\ auto\isanewline
\ \ \ \ \isacommand{have}\isamarkupfalse%
\ {\isachardoublequoteopen}a\ {\isacharasterisk}{\kern0pt}\ c\ {\isacharequal}{\kern0pt}\ a\ {\isacharasterisk}{\kern0pt}\ {\isacharparenleft}{\kern0pt}Suc\ nat{\isacharparenright}{\kern0pt}{\isachardoublequoteclose}\isanewline
\ \ \ \ \ \ \isacommand{apply}\isamarkupfalse%
{\isacharparenleft}{\kern0pt}subst\ \ Suc{\isacharparenright}{\kern0pt}\isanewline
\ \ \ \ \ \ \isacommand{{\isachardot}{\kern0pt}{\isachardot}{\kern0pt}}\isamarkupfalse%
\isanewline
\ \ \ \ \isacommand{hence}\isamarkupfalse%
\ {\isadigit{4}}{\isacharcolon}{\kern0pt}\ {\isachardoublequoteopen}a\ {\isacharasterisk}{\kern0pt}\ c\ {\isacharequal}{\kern0pt}\ a\ {\isacharplus}{\kern0pt}\ a\ {\isacharasterisk}{\kern0pt}\ nat{\isachardoublequoteclose}\isanewline
\ \ \ \ \ \ \isacommand{by}\isamarkupfalse%
{\isacharparenleft}{\kern0pt}subst\ {\isacharparenleft}{\kern0pt}asm{\isacharparenright}{\kern0pt}\ mult{\isacharunderscore}{\kern0pt}Suc{\isacharunderscore}{\kern0pt}right{\isacharparenright}{\kern0pt}\isanewline
\isanewline
\ \ \ \ \isacommand{hence}\isamarkupfalse%
\ {\isadigit{5}}{\isacharcolon}{\kern0pt}\ {\isachardoublequoteopen}b\ {\isacharasterisk}{\kern0pt}\ Suc\ d\ {\isacharequal}{\kern0pt}\ b\ {\isacharplus}{\kern0pt}\ b\ {\isacharasterisk}{\kern0pt}\ d{\isachardoublequoteclose}\isanewline
\ \ \ \ \ \ \isacommand{apply}\isamarkupfalse%
\ {\isacharparenleft}{\kern0pt}subst\ mult{\isacharunderscore}{\kern0pt}Suc{\isacharunderscore}{\kern0pt}right{\isacharparenright}{\kern0pt}\isanewline
\ \ \ \ \ \ \isacommand{{\isachardot}{\kern0pt}{\isachardot}{\kern0pt}}\isamarkupfalse%
\isanewline
\ \ \ \ \isacommand{have}\isamarkupfalse%
\ {\isachardoublequoteopen}a\ {\isacharplus}{\kern0pt}\ a\ {\isacharasterisk}{\kern0pt}\ nat\ {\isasymle}\ b\ {\isacharplus}{\kern0pt}\ b\ {\isacharasterisk}{\kern0pt}\ d{\isachardoublequoteclose}\isanewline
\ \ \ \ \ \ \isacommand{by}\isamarkupfalse%
{\isacharparenleft}{\kern0pt}rule\ add{\isacharunderscore}{\kern0pt}mono{\isacharbrackleft}{\kern0pt}OF\ {\isacartoucheopen}a\ {\isasymle}\ b{\isacartoucheclose}\ {\isadigit{3}}{\isacharbrackright}{\kern0pt}{\isacharparenright}{\kern0pt}\isanewline
\ \ \ \ \isacommand{then}\isamarkupfalse%
\ \isacommand{have}\isamarkupfalse%
\ {\isachardoublequoteopen}a\ {\isacharasterisk}{\kern0pt}\ c\ {\isasymle}\ b\ {\isacharplus}{\kern0pt}\ b\ {\isacharasterisk}{\kern0pt}\ d{\isachardoublequoteclose}\isanewline
\ \ \ \ \ \ \isacommand{apply}\isamarkupfalse%
{\isacharparenleft}{\kern0pt}subst\ {\isadigit{4}}{\isacharparenright}{\kern0pt}\isanewline
\ \ \ \ \ \ \isacommand{{\isachardot}{\kern0pt}}\isamarkupfalse%
\isanewline
\ \ \ \ \isacommand{then}\isamarkupfalse%
\ \isacommand{show}\isamarkupfalse%
\ {\isacharquery}{\kern0pt}thesis\isanewline
\ \ \ \ \ \ \isacommand{apply}\isamarkupfalse%
{\isacharparenleft}{\kern0pt}subst\ {\isadigit{5}}{\isacharparenright}{\kern0pt}\isanewline
\ \ \ \ \ \ \isacommand{{\isachardot}{\kern0pt}}\isamarkupfalse%
\isanewline
\ \ \isacommand{qed}\isamarkupfalse%
\isanewline
\isacommand{qed}\isamarkupfalse%
%
\endisatagproof
{\isafoldproof}%
%
\isadelimproof
%
\endisadelimproof
\ \isanewline
\isanewline
\isanewline
\isanewline
%
\isadelimtheory
\isanewline
%
\endisadelimtheory
%
\isatagtheory
\isacommand{end}\isamarkupfalse%
%
\endisatagtheory
{\isafoldtheory}%
%
\isadelimtheory
%
\endisadelimtheory
%
\end{isabellebody}%
\endinput
%:%file=Week_1_sol.tex%:%
%:%6=1%:%
%:%11=2%:%
%:%12=2%:%
%:%13=3%:%
%:%14=4%:%
%:%23=6%:%
%:%27=8%:%
%:%28=9%:%
%:%29=10%:%
%:%30=11%:%
%:%31=12%:%
%:%32=13%:%
%:%36=16%:%
%:%40=20%:%
%:%41=21%:%
%:%42=22%:%
%:%44=39%:%
%:%45=39%:%
%:%47=41%:%
%:%48=42%:%
%:%49=43%:%
%:%50=44%:%
%:%51=45%:%
%:%52=46%:%
%:%53=47%:%
%:%54=48%:%
%:%55=49%:%
%:%56=50%:%
%:%57=51%:%
%:%65=54%:%
%:%66=54%:%
%:%67=55%:%
%:%68=55%:%
%:%69=56%:%
%:%70=56%:%
%:%71=57%:%
%:%72=57%:%
%:%73=58%:%
%:%74=58%:%
%:%75=59%:%
%:%76=59%:%
%:%77=60%:%
%:%78=60%:%
%:%79=61%:%
%:%80=61%:%
%:%81=62%:%
%:%82=63%:%
%:%83=63%:%
%:%84=64%:%
%:%85=64%:%
%:%86=65%:%
%:%87=65%:%
%:%88=66%:%
%:%89=66%:%
%:%90=66%:%
%:%91=67%:%
%:%92=67%:%
%:%93=68%:%
%:%94=68%:%
%:%95=69%:%
%:%96=69%:%
%:%97=70%:%
%:%98=71%:%
%:%99=71%:%
%:%100=72%:%
%:%101=72%:%
%:%102=73%:%
%:%103=73%:%
%:%104=74%:%
%:%105=75%:%
%:%106=75%:%
%:%107=76%:%
%:%108=76%:%
%:%109=77%:%
%:%110=77%:%
%:%111=78%:%
%:%112=79%:%
%:%113=79%:%
%:%114=80%:%
%:%115=81%:%
%:%116=81%:%
%:%117=82%:%
%:%118=83%:%
%:%119=83%:%
%:%120=84%:%
%:%121=84%:%
%:%122=85%:%
%:%123=86%:%
%:%124=86%:%
%:%125=87%:%
%:%126=87%:%
%:%127=88%:%
%:%128=89%:%
%:%129=89%:%
%:%130=90%:%
%:%131=90%:%
%:%132=91%:%
%:%133=92%:%
%:%143=99%:%
%:%147=101%:%
%:%148=102%:%
%:%149=103%:%
%:%151=126%:%
%:%152=126%:%
%:%153=127%:%
%:%154=127%:%
%:%155=127%:%
%:%156=128%:%
%:%157=129%:%
%:%159=132%:%
%:%160=133%:%
%:%161=134%:%
%:%162=135%:%
%:%163=136%:%
%:%164=137%:%
%:%165=138%:%
%:%166=139%:%
%:%167=140%:%
%:%168=141%:%
%:%169=142%:%
%:%170=143%:%
%:%171=144%:%
%:%172=145%:%
%:%173=146%:%
%:%174=147%:%
%:%175=148%:%
%:%176=149%:%
%:%177=150%:%
%:%178=151%:%
%:%179=152%:%
%:%180=153%:%
%:%181=154%:%
%:%181=155%:%
%:%182=156%:%
%:%183=157%:%
%:%184=158%:%
%:%185=159%:%
%:%189=165%:%
%:%193=165%:%
%:%194=165%:%
%:%195=166%:%
%:%196=166%:%
%:%197=167%:%
%:%198=167%:%
%:%199=168%:%
%:%200=168%:%
%:%201=168%:%
%:%202=169%:%
%:%203=169%:%
%:%204=170%:%
%:%205=170%:%
%:%206=171%:%
%:%207=171%:%
%:%208=172%:%
%:%209=172%:%
%:%210=172%:%
%:%211=173%:%
%:%212=173%:%
%:%213=174%:%
%:%214=174%:%
%:%215=175%:%
%:%216=175%:%
%:%217=176%:%
%:%218=176%:%
%:%219=177%:%
%:%220=177%:%
%:%221=178%:%
%:%222=179%:%
%:%223=180%:%
%:%224=181%:%
%:%225=181%:%
%:%226=182%:%
%:%227=183%:%
%:%228=183%:%
%:%229=184%:%
%:%230=184%:%
%:%231=185%:%
%:%232=186%:%
%:%233=187%:%
%:%234=188%:%
%:%235=188%:%
%:%236=189%:%
%:%237=189%:%
%:%238=190%:%
%:%239=190%:%
%:%240=191%:%
%:%241=191%:%
%:%242=192%:%
%:%243=192%:%
%:%244=193%:%
%:%245=193%:%
%:%246=194%:%
%:%247=194%:%
%:%248=195%:%
%:%249=195%:%
%:%250=196%:%
%:%251=196%:%
%:%252=197%:%
%:%253=197%:%
%:%254=198%:%
%:%255=198%:%
%:%256=199%:%
%:%257=199%:%
%:%258=200%:%
%:%259=200%:%
%:%260=201%:%
%:%261=201%:%
%:%262=202%:%
%:%263=202%:%
%:%264=203%:%
%:%265=203%:%
%:%266=204%:%
%:%267=204%:%
%:%268=205%:%
%:%269=205%:%
%:%270=206%:%
%:%271=206%:%
%:%272=207%:%
%:%273=207%:%
%:%274=208%:%
%:%275=208%:%
%:%276=209%:%
%:%277=209%:%
%:%278=210%:%
%:%279=210%:%
%:%280=211%:%
%:%281=211%:%
%:%282=212%:%
%:%283=212%:%
%:%284=213%:%
%:%285=214%:%
%:%286=214%:%
%:%287=215%:%
%:%288=215%:%
%:%289=216%:%
%:%290=216%:%
%:%291=217%:%
%:%292=217%:%
%:%293=218%:%
%:%294=218%:%
%:%295=219%:%
%:%296=219%:%
%:%297=219%:%
%:%298=220%:%
%:%299=220%:%
%:%300=221%:%
%:%301=221%:%
%:%302=222%:%
%:%303=222%:%
%:%304=222%:%
%:%305=223%:%
%:%306=223%:%
%:%307=224%:%
%:%308=224%:%
%:%309=225%:%
%:%310=225%:%
%:%311=226%:%
%:%319=226%:%
%:%320=227%:%
%:%321=228%:%
%:%322=229%:%
%:%325=230%:%
%:%330=231%:%
